\section{Intrinsic fibrations and coefficient localizers}
\subsection{Truncation and the boundary homomorphism}
\begin{claim}{D.1}{Known}{The proper-localizer criterion, including its local form}
For $n\geq0$, put $W_n=\{u:\tau_{\leq n}Bu\text{ is an equivalence}\}$. A functor $p:E\to B$ belongs to $\Fib_{W_n}$ if and only if every strict pullback square along $X\to B$ becomes a homotopy pullback after applying $\tau_{\leq n}B(-)$. Equivalently, it suffices to require this for all $X$ having a terminal object. Another equivalent test is that, for every triangle $X\to Y\to B$ with $X$ having a terminal object and $BY$ $n$-connected, the functor $E\times_BX\to E\times_BY$ belongs to $W_n$. For $p\in\Liss_{W_n}$ or $p\in\Prop_{W_n}$, it suffices to test the squares along the objects $b:*\to B$.
\end{claim}
\begin{proof}
These are, respectively, Proposition 6.2.3(i)--(iv) and Proposition 6.2.4 of \citet{RESCisinski2006}. Their hypothesis is that the fundamental localizer is proper. Section 9.2.1 of that source establishes properness of $W_n$, and Corollary 9.2.15 identifies its maps with the stated homotopy-group test. The associated localization is the category of $n$-types; its homotopy pullbacks are the pullbacks appearing here. This gives a criterion for arbitrary functors, without silently replacing strict pullbacks by homotopy pullbacks in spaces. In particular,
\[
\Fib_{W_n}\cap\Liss_{W_n}
 =\{p\in\Liss_{W_n}:\tau_{\leq n}BE_b\simeq
 \operatorname{fib}_b(\tau_{\leq n}BE\to\tau_{\leq n}BB)\},
\]
with the analogous equality for $\Prop_{W_n}$. The equivalences in this formula are the canonical fibre comparison maps, for every $b$; abstract equivalence of their endpoints is insufficient.
\end{proof}

\begin{claim}{D.2}{Proved}{Lead 1 is true}
Suppose $p:E\to B$ is a small-category functor such that, for every $X\to B$, the canonical map $B(E\times_BX)\to BE\times^h_{BB}BX$ is a weak equivalence. For $n\geq1$, it belongs to $\Fib_{W_n}$ if and only if every connecting homomorphism
\[
\partial_b:\pi_{n+1}(BB,b)\longrightarrow\pi_n(BE_b,x)
\]
is zero, for every $b$ and every vertex $x\in E_b$. For $n=0$, the criterion is that the monodromy action of $\pi_1(BB,b)$ on $\pi_0(BE_b)$ is trivial. Empty fibres and components outside the image of $BE$ are allowed.
\end{claim}
\begin{proof}
Let $F=BE_b$, choose $x\in F$, and replace $BE\to BB$ by a Serre fibration to write its long exact sequence. This replacement is used only to compute its homotopy fibre; the hypothesis identifies that fibre with the specified strict fibre. For $n\geq1$, the canonical map
\[
\tau_{\leq n}F\longrightarrow
\operatorname{fib}_b(\tau_{\leq n}BE\to\tau_{\leq n}BB)
\tag{D.2.1}
\]
is an isomorphism on all homotopy groups below $n$, including the component sets and the actions in degree zero. In degree $n$ its map is the quotient
$\pi_nF\to\pi_nF/\operatorname{im}\partial_b$.
For $n=1$ the image is the normal subgroup specified by the exact homotopy sequence, so this statement is a statement about groups, not about abelianization. Thus (D.2.1) is an equivalence if and only if $\partial_b=0$. Necessity follows at once from D.1 applied to $X=*$. In degree zero, the fibre of $\pi_0BE\to\pi_0BB$ over a component is the orbit set of $\pi_0F$ under $\pi_1BB$. The quotient map to this orbit set is bijective exactly when every orbit is a singleton. This proves necessity for $n=0$.

For sufficiency let $v:X\to B$ be arbitrary and write $E_X=E\times_BX$. For every $x'\in X$, its homotopy fibre is $BE_{v(x')}$. Naturality of the connecting homomorphism identifies the boundary for $E_X\to X$ with the composite of the boundary for $E\to B$ and $\pi_{n+1}(Bv)$. It is therefore zero. Apply (D.2.1) to both fibrations and compare the fibres of
\[
\tau_{\leq n}BE_X\longrightarrow
\tau_{\leq n}BE\times^h_{\tau_{\leq n}BB}\tau_{\leq n}BX
\]
over $\tau_{\leq n}BX$. These fibre maps are equivalences. A map of spaces over a fixed base that is an equivalence on every homotopy fibre is an equivalence, proving the square criterion in D.1. For $n=0$, monodromy after pullback is the restriction of the original action, so it remains trivial and the same argument works with component sets. On a component with empty homotopy fibre, the total pullback is empty on both sides. Every component meeting a nonempty fibre is treated with one basepoint in each of its fibre components. No surjectivity assumption on $\pi_0BE\to\pi_0BB$ was used.
\end{proof}

\begin{claim}{D.3}{Proved}{An elementary complete criterion at $W_0$}
For an arbitrary functor $p:E\to B$, put $E_j=E\times_BB_j$ for each connected component $B_j$ of $B$. Then $p\in\Fib_{W_0}$ precisely when (i) every $\pi_0(E_b)\to\pi_0(E_j)$ is bijective, and (ii) every arrow $\beta:b\to b'$ in $B_j$ has a lift in every prescribed component of $E_j$.
\end{claim}
\begin{proof}
Pull back the $W_0$-map $*\to B_j$ to obtain (i). Pull back the two inclusions $*\to[1]$ along the base change of $p$ by $\beta:[1]\to B$. By (i), each prescribed component has exactly one connected fibre over either endpoint. These fibres must be connected to each other in the total pullback. A zigzag connecting the layers must contain an arrow over $\beta$, proving (ii). Conversely, over a connected component $X_a$ of any $X\to B_j$, label objects of $E\times_BX_a$ by their component in $E_j$. Condition (i) says each labelled fibre is nonempty and connected. Condition (ii) connects these fibres along each arrow of $X_a$. A zigzag in $X_a$ therefore connects precisely the objects with the same label. This proves, naturally in $X$,
$\pi_0(E\times_BX)=\pi_0E\times_{\pi_0B}\pi_0X$.
This formula preserves every $W_0$-map and proves the result, including empty categories.
\end{proof}

\begin{claim}{D.4}{Refuted}{Intrinsic fibrations need not be a right lifting class}
Let $W$ be a fundamental localizer with $W_\infty\subseteq W\subseteq W_0$. Then $\Fib_W$ is not the class of maps with the right lifting property against any class of functors, and in particular against no set. This applies to every $W_n$, $W_\infty$, nonzero ordinary homology, and the nontrivial coefficient localizers considered here.
\end{claim}
\begin{proof}
For each integer $r\geq1$, let $C_r$ be the fence poset on $0,\ldots,r$, with a strict order relation between adjacent vertices directed from the even vertex to the odd vertex and no other strict relations. Its nerve is an interval. Let $A_r=C_r^{\triangleright}$, obtained by adjoining a terminal object $\infty_r$, and let $p_r:A_r\to[1]$ take $C_r$ to $0$ and $\infty_r$ to $1$.
For any $X\to[1]$, the pullback is the Grothendieck construction of the diagram on $X$ that has fibre $C_r$ at colour $0$ and $*$ at colour $1$. The map of this diagram to the constant point diagram is objectwise a weak equivalence. The homotopy-colimit comparison for the Grothendieck construction gives $A_r\times_{[1]}X\to X$ in $W_\infty$. Thus every base change of $p_r$ is a weak equivalence, and two-out-of-three shows $p_r\in\Fib_W$.

Form the countable product $p:\prod_rA_r\to\prod_r[1]$. Both its domain and codomain have terminal objects. The inclusion of the all-zero object $*\to\prod_r[1]$ belongs to $W_\infty$. Its pullback along $p$ is $\prod_rC_r\to\prod_rA_r$. The first category is disconnected: the objects $(0)_r$ and $(r)_r$ cannot be joined by a finite zigzag. A zigzag with $k$ arrows changes each coordinate by at most $k$, whereas their $r$th coordinates differ by $r$. Choosing $r>k$ is a contradiction. The second category is connected, so this pullback is not in $W_0$, hence not in $W$.

Any right lifting class is closed under products: solve a lifting square in each coordinate and take the product of the lifts. This contradicts the displayed family. All categories and compositions have been specified, and the program checks the individual fences and their distance bound; disconnectedness of the infinite product is proved by the finite-zigzag argument, not inferred from finite computations.
\end{proof}

\subsection{Locally constant derived coefficients}
\begin{claim}{D.5}{Known}{The coefficient tests and their variance}
For a nonzero commutative ring $k$, a functor $u:A\to B$ lies in $W_{D(k)}^{\mathrm{lc}}$ exactly when it is bijective on components, induces isomorphisms on $\pi_1$ at all basepoints, and induces $k$-homology equivalences of the simply connected universal covers of corresponding components. Consequently $W_{D(\mathbb Z)}^{\mathrm{lc}}=W_\infty$. For $k=\mathbb Q$ the higher test is $\pi_i(Bu)\otimes\mathbb Q$ being an isomorphism for every $i\geq2$. For sets and groupoids as coherent coefficients, the locally constant classes are $W_1$ and $W_2$, respectively.
\end{claim}
\begin{proof}
Evaluation at a basepoint of a connected space $Z$ identifies $\Fun(Z,D(k))$ with modules over the specified loop algebra $C_*(\Omega Z;k)$. Evaluation is conservative and preserves colimits; its left adjoint has monad $C_*(\Omega Z;k)\otimes_k-$. The monadic comparison follows from the bar resolution, or \citet[Theorem 4.7.3.5]{GHIJLurie2017}. Under this comparison restriction along a based map is restriction of scalars along its loop-algebra map. Restriction is an equivalence only if the unit on the regular source module is an equivalence, hence only if that algebra map is a quasi-isomorphism; conversely a quasi-isomorphism induces an equivalence of module categories by the two-sided bar resolution.
The degree-zero map is $k[\pi_1A]\to k[\pi_1B]$. Since this is the specified map on bases, it is invertible only if the group map is bijective. Each loop-space component then reduces to the identity component, the loop space of the universal cover. The bar and cobar constructions on simply connected spaces show that a map of their loop chains is a $k$-homology equivalence if and only if the original map is one: use the tensor-length filtration; after desuspending reduced chains all generators have positive degree, so the filtration is finite in each total degree. Tensoring bounded-below free chain complexes preserves the requisite quasi-isomorphisms. This gives the asserted universal-cover criterion. With $k=\mathbb Z$, homological Whitehead on the simply connected covers proves the equality with $W_\infty$. The rational identification is also precisely \citet[Theorem 3]{RRivera2021}.
For $r$-truncated coefficients, restriction depends only on $\tau_{\leq r+1}Z$: the cores of all diagram categories $\Fun([m],\Spaces_{\leq r})$ are $(r+1)$-types. Conversely, evaluating the unit on the free object at a basepoint detects $\tau_{\leq r}\Omega Z$. This proves $W^{\mathrm{lc}}_{\Spaces_{\leq r}}=W_{r+1}$. The zero ring is excluded because then all coefficient categories are terminal.
\end{proof}

\begin{claim}{D.6}{Proved}{Rational local coefficients require only the degree-two boundary}
Under the homotopy-pullback hypothesis of D.2, $p:E\to B$ belongs to $\Fib_{W_{D(\mathbb Q)}^{\mathrm{lc}}}$ if and only if every boundary $\pi_2(BB,b)\to\pi_1(BE_b,x)$ is zero. There is no vanishing condition on the higher rational boundaries. In particular the stronger proposed test in Part D(3) is false.
\end{claim}
\begin{proof}
Assume these boundaries vanish. For a $W_{D(\mathbb Q)}^{\mathrm{lc}}$-map $u:X\to Y$ over any base change of $B$, compare the homotopy fibration sequences with common fibre $F$. The isomorphism on fundamental groups of $X\to Y$ identifies the monodromy orbit sets, hence the components of the two total spaces. On each such component the vanishing boundary gives an injection $\pi_1F\to\pi_1E_X$ and an exact extension by the stabilizer subgroup of $\pi_1X$ acting on the chosen fibre component. The two extensions have isomorphic kernels and quotients, so the total fundamental groups are isomorphic. Tensor the higher long exact sequences with $\mathbb Q$. For degrees at least two all groups in the relevant five-term segments are abelian, and the zero degree-two boundary permits the sequence to end with $\pi_2E_X\to\pi_2X\to0$. Exactness of rational tensoring and the five lemma give isomorphisms on every higher rational homotopy group. D.5 proves sufficiency.

Conversely take $z\in\pi_2(BB,b)$ with boundary $a\ne1$ in $\pi_1F$. Pull back the fibration to a based sphere representing $z$, and compare its pullback along a degree-$m$ sphere map. A degree-$m$ map of $S^2$, for $m\ne0$, belongs to $W_{D(\mathbb Q)}^{\mathrm{lc}}$. On the chosen total-space components the induced fundamental-group map is
\[
 \pi_1F/\langle a^m\rangle\longrightarrow\pi_1F/\langle a\rangle.
\]
Both displayed cyclic subgroups are normal, by the exact sequence of the pulled-back fibration over $S^2$. If $a$ has finite order $d>1$, choose $m=d$; if it has infinite order, choose $m=2$. The map then has a nontrivial kernel and is not a coefficient equivalence. We justify the strict categorical witness. The weak equivalence $NB\to\operatorname{Ex}^{\infty}NB$, together with right properness of simplicial sets, gives a Quillen equivalence of the slice model categories. Thus the sphere triangle in spaces over $BB$ has a representative $K'\to K\to NB$ in simplicial sets: represent the target by a fibrant object of the slice and then represent the map to it. Apply the category-of-simplices construction to this strict triangle. Last-vertex evaluation gives functors $\Delta/K'\to\Delta/K\to B$, and the natural weak equivalences $B(\Delta/K)\simeq|K|$ identify their realized triangle with the chosen triangle in the slice of spaces. Since its first map is a rational local-coefficient equivalence, so is $\Delta/K'\to\Delta/K$. The homotopy-pullback hypothesis on $p$ identifies the strict categorical pullbacks of this triangle with the homotopy pullbacks used above. This transfers the nontrivial fundamental-group kernel to a categorical witness. Thus intrinsic membership forces $a=1$.

To exhibit the unnecessary higher condition, take a categorical homotopy-pullback representative of the path fibration over $S^3$. For example factor $*\to K$, where $K$ is the face poset of the boundary of a $4$-simplex, into a Thomason weak equivalence followed by a Thomason fibration. Right properness makes its strict pullbacks homotopy pullbacks. Its total category is weakly contractible and its fibre has type $\Omega S^3$. Its degree-two boundary is zero, while its boundary $\pi_3S^3\to\pi_2\Omega S^3$ is an isomorphism $\mathbb Z\to\mathbb Z$. The preceding equivalence places this map in the stated intrinsic class. For definiteness the factorization is constructed as follows. Starting from $P_0=*$ over $K$, form $P_{r+1}$ by attaching one copy of $c\operatorname{Sd}^2\Delta^m$ for every commutative lifting square from $c\operatorname{Sd}^2\Lambda_k^m$ into $P_r$ over $K$, for all $m\geq1$ and $0\leq k\leq m$, and set $P=\colim_{r<\omega}P_r$. Each attachment is the categorical pushout: freely adjoin its objects and arrows and impose precisely its categorical relations and the relations of the attaching functor. This specifies all objects, arrows, and composition of $P$ as equivalence classes of finite words. Finite presentability of the attaching domains gives the right lifting property; Thomason's generating-trivial-cofibration theorem gives $*\to P\in W_\infty$. Thus this is a completely specified infinite-category example, not a finite category certificate.
\end{proof}

\begin{claim}{D.7}{Open}{The remaining coefficient classification}
For arbitrary functors without the hypothesis of D.2, a comparable explicit fibre criterion for $\Fib_{W_{D(\mathbb Q)}^{\mathrm{lc}}}$ is not established here. Nor do we establish a complete characterization of $\Fib_{W_{D(\mathbb F_p)}^{\mathrm{lc}}}$, even among the homotopy-pullback functors of D.2. The integral, set, and groupoid cases reduce exactly to $W_\infty,W_1,W_2$ by D.5. Merely declaring higher boundary maps zero after tensoring with a field is not a proof of the unresolved mod-$p$ case: integral fundamental groups are retained by that coefficient localizer.
\end{claim}

\subsection{Homology and covering maps}
\begin{claim}{D.8}{Refuted}{A connected intrinsic homology covering need not be trivial}
Let $p$ be a prime. Any categorical covering whose monodromy factors through a finite $p$-group belongs to $\Fib_{W_{H\mathbb F_p}}$. In particular the regular covering $EC_p\to BC_p$ is a connected nontrivial covering in this class. Thus the proposed assertion for every nonzero spectrum in Part D(4) is false.
\end{claim}
\begin{proof}
Write $G$ for the finite $p$-group and $S$ for the fibre set. The permutation module $V=\mathbb F_p[S]$ has a finite filtration with trivial $G$-action on its successive quotients. Here is a proof independent of a semisimplicity assertion. The augmentation ideal $J$ of $\mathbb F_p[G]$ is nilpotent. Induct on $|G|$: choose a central subgroup $C_p\subseteq G$, let $z$ generate it, and use $(z-1)^p=0$. The quotient algebra by $(z-1)$ is $\mathbb F_p[G/C_p]$. If its augmentation ideal has nilpotence exponent $d$, then $J^d\subseteq(z-1)$ and $J^{dp}=0$. The filtration $V\supseteq JV\supseteq\cdots$ is therefore finite, also for infinite $S$, and $G$ acts trivially on every quotient.

After any base change $X\to B$, the covering is classified by the restricted permutation action. Its homology is the homology of $X$ with this local coefficient module. A mod-$p$ homology equivalence $Y\to X$ induces an isomorphism on homology with any constant $\mathbb F_p$-vector-space coefficient, since tensoring over a field is exact. Apply the long exact sequences of homology with local coefficients successively to the finite filtration. They prove the same assertion with coefficient $V$, hence for the pulled-back coverings. This proves the universal intrinsic condition, component by component.
For the concrete finite example, $BC_p$ has one object and arrows $\mathbb Z/p$; $EC_p$ has objects $\mathbb Z/p$ and one arrow $s\to t$, mapped to $t-s$. Composition is addition of differences. Every base arrow has a unique lift from any source and to any target. The total category is indiscrete and contractible; the fibre has $p$ objects. The program verifies all category laws and both unique-lift conditions for $p=2,3,5,7$, including the nilpotence matrix identity. The proof applies to every prime.
\end{proof}

\begin{claim}{D.9}{Open}{The remaining homology classification}
A complete classification for arbitrary nonzero spectra $E$, arbitrary coverings, and all $W_\infty$-smooth and proper functors is not proved here. D.8 shows that the answer must depend on the coefficient theory and on monodromy. D.12 below gives complete covering classifications for ordinary integral and rational homology. Product projections and unions of components give intrinsic examples. Neither those results nor boundary homomorphisms alone give the classification for every spectrum. No finite computation in this return claims to test all base changes in the definition of $\Fib_{W_E}$.
\end{claim}

\subsection{The eight inherited finite maps}
\begin{claim}{D.10}{Proved}{BISITE B.15 has no changed entries}
For the maps defined below, the entries are as follows; $Q\mathrm{lc}$ denotes $W_{D(\mathbb Q)}^{\mathrm{lc}}$.
\begin{center}\begin{tabular}{lcccc}
\toprule Map & Smooth, both & Proper, both & Intrinsic $W_1$ & Intrinsic $Q\mathrm{lc}$\\\midrule
D.3: $p:E\to B$ & yes&yes&yes&yes\\
D.3: $q:B\to[1]$ & no&yes&no&no\\
D.3: $qp$ & no&yes&no&no\\
D.4: $p:E\to B$ & yes&yes&no&no\\
D.4: $q:B\to[1]$ & no&yes&yes&no\\
D.4: $qp$ & no&yes&no&no\\
D.5: $H\to K$ & yes&yes&no&no\\
D.5: $q:R\to[1]$ & no&yes&no&no\\\bottomrule
\end{tabular}\end{center}
\end{claim}
\begin{proof}
We specify all categories. For D.3, $B$ has objects $a,b$, endomorphisms $1_a$ and $C_2$ at $b$, and two arrows $a\to b$ labelled by $C_2$, with postcomposition given by addition. Put $L(a)=L(b)=C_2$, and let a labelled arrow act by addition. Set $E=\int_BL$ and let $q$ record the two colours. Then $p$ is a covering. Coverings preserve the two classes: the base map identifies fundamental groups and their permutation actions, while the universal covers of corresponding total components identify with the universal covers of the base components. D.5 proves preservation of both tests after every base change. The incoming lift categories for $q$ and $qp$ include two-point discrete categories; their outgoing lifts have initial objects. Pulling $\{0\}\to[1]$ back along $q$ gives $*\to B\simeq BC_2$, failing on $\pi_1$. Along $qp$ it gives a discrete two-point fibre into the connected contractible $E$, failing on $\pi_0$.

Let $K$ have objects the nonempty proper subsets of $\{0,1,2,3\}$, with inclusion arrows. For $x\subseteq y\subseteq z$ let $c(x,y,z)\in C_2$ be $1$ only when $(x,y,z)=(\{0\},\{0,1\},\{0,1,2\})$. Let $H$ have the objects of $K$ and arrows $(x\subseteq y,m)$, $m\in C_2$, with
$(y\subseteq z,n)(x\subseteq y,m)=(x\subseteq z,m+n+c(x,y,z))$.
There are no three strict composable inclusions, so the normalized cocycle identity holds. The projection is a bifibration, by unique solution for the group label of every cartesian or cocartesian factorization. Its fibre is $BC_2$. The chosen cocycle evaluates nontrivially on the triangulated sphere $BK$, so the connecting map $\pi_2BK\to C_2$ is onto. Equivalently, collapse a spanning tree in the edge presentation: all group generators become one central involution; the triangle relations outside the marked triangle form a disk and kill all horizontal generators; the marked triangle kills the involution. Thus $\pi_1BH=0$. In $H\times_KH$ the same presentation leaves $C_2^2/\langle(1,1)\rangle=C_2$. The accompanying program constructs all arrows and relations and performs explicit Tietze eliminations to these groups. Thus $*\to K$ tests failure for $W_1$, and $H\to K$ itself tests failure for rational local coefficients: it is a rational higher-homotopy equivalence with trivial fundamental groups, but its self-pullback has fundamental group $C_2$. This also proves that the nonzero boundary in D.6 is detected integrally.

For D.4 adjoin $a$ to $K$, with $B(a,x)=K(\{0\},x)$, no arrows into $a$ except its identity, and the forced composition. Let $r:B\to K$ take $a$ to $\{0\}$ and otherwise be the identity; set $E=B\times_KH$. The inclusion $K\to B$ is right adjoint to $r$, and the same adjunction pulls back to $H\to E$, so $BB\simeq S^2$ and $BE\simeq BH$. The map $p$ is a bifibration and inherits both failures from $H\to K$. The maps $q$ and $qp$ are opfibrations. An object of $K$ not containing $0$ has no incoming lift from colour $0$, proving nonsmoothness. The fibre of $qp$ at $0$ is $BC_2$, and its inclusion into $E$ fails on $\pi_1$. The map $q$ is universally a $W_1$-equivalence: its fibre diagram consists of a point and the simply connected $K$, and groupoid completion of its Grothendieck construction contracts each fibre to its chosen section object. The unique paths in the completed fibre make this contraction natural over any base change. Two-out-of-three gives intrinsic $W_1$-membership. Its fibre inclusion $*\to B\simeq S^2$ is not a rational higher-homotopy equivalence, excluding the other membership.

Finally $R$ has $b,a$, endomorphism group $C_2$ at $b$, identity at $a$, and a unique arrow $b\to a$ fixed by precomposition with $C_2$. It has terminal object $a$, is an opfibration over $[1]$, and has incoming lift category $BC_2$ at $a$. The inclusion of its colour-zero fibre fails on fundamental groups. These checks prove all table entries. The independently written program verifies functoriality, every outgoing and incoming lifting category, and the fundamental-group presentations; it does not infer universal preservation merely from a finite list of pullbacks.
\end{proof}

\subsection{The axioms of XVI, section 4}
\begin{claim}{D.11}{Proved}{The requested fibration-axiom table}
Let $\mathcal F_W$ denote Grothendieck fibrations with $W$-equivalent transport functors. We use the printed meanings: (7a) is $\mathcal F_W\subseteq\Fib_W$; (7b) says a $W$-map between two members of $\mathcal F_W$ over a common base is universally in $W$; (7c) says $\mathcal F_W\cap W$ is universally in $W$; and (7F) is existence of $W$-then-$\Fib_W$ factorizations. Then
\begin{center}\begin{tabular}{lccccc}
\toprule $W$ & (7a)&(7b)&(7c)&(7F)&$\Liss_W\cap\Prop_W\subseteq\Fib_W$\\\midrule
$W_0$ & no&no&no&yes&no\\
$W_1$ & no&no&no&yes&no\\
$W_n$, $n\geq2$ & no&no&no&yes&no\\
$W_{\mathbb Q}$ & no&no&no&no&no\\
$W_{D(\mathbb Q)}^{\mathrm{lc}}$ & no&no&no&no&no\\
$W_\infty$ & yes&yes&yes&yes&yes\\\bottomrule
\end{tabular}\end{center}
\end{claim}
\begin{proof}
For $W_0$, the regular $C_2$-covering has connected base and total category and a disconnected fibre. It is a bifibration, is itself in $W_0$, and is not universally in $W_0$. For $W_{\mathbb Q}$ the same covering is a rational homology equivalence, since averaging contracts the positive rational group homology of $C_2$. The same fibre detects failure. For $W_1$ and $W_{D(\mathbb Q)}^{\mathrm{lc}}$, the finite bifibration $H\to K$ of D.10 is itself in the respective class but not universally in it. These examples each refute (7c) and (7b), the latter by comparison with $\id_B$ over $B$, as well as (7a) and the last column.

For general finite $n$, let $K_n$ be the face poset of the boundary of an $(n+2)$-simplex, so $BK_n\simeq S^{n+1}$. The smooth and proper categorical path replacement established in Part C has contractible total category and fibre $\Omega S^{n+1}$. Its projection and $*\to K_n$ belong to $W_n$, but its fibre is not $n$-connected. This proves the last-column failure. To obtain the required Grothendieck fibration for (7c), start with a Thomason-fibration replacement $P\to K_n$ of $*\to K_n$ and form the contravariant comma diagram $b\mapsto b\backslash P$. Its Grothendieck construction $Q\to K_n$ is a Grothendieck fibration. The forgetful functor $Q\to P$ has the right adjoint $e\mapsto(p(e),e,1)$ and hence is a weak equivalence. Right properness identifies each fibre $b\backslash P$ with the homotopy fibre, and identifies all its transports as weak equivalences. Thus $Q\to K_n$ lies in $\mathcal F_{W_n}\cap W_n$ and has the same non-$W_n$ fibre. This refutes (7c). Comparison with $\id_{K_n}$ refutes (7b). Pulling the $W_n$-map $*\to K_n$ back refutes (7a).

For $W_\infty$, Quillen's Theorem B applied to a Grothendieck fibration with weakly equivalent transports identifies all strict base changes as homotopy pullbacks. This proves (7a). A weak equivalence between the total spaces of two such homotopy fibrations over the same base is a weak equivalence on all homotopy fibres; after any further base change it remains so, proving (7b) and (7c). Smoothness together with properness gives the same homotopy-pullback property, by the fibre-to-comma comparisons on both sides of an arrow. Thomason's right proper model structure supplies (7F). For $W_n$, properness in \citet[Section 9.2.1]{RESCisinski2006}, together with Theorem 5.2.15 and Example 6.1.14 of that source, supplies (7F).

We justify the two negative (7F) entries. For an accessible fundamental localizer $W$, existence of the stated factorizations is equivalent to properness of $W$. For the less immediate direction, let $W_{\rm asph}$ be the fundamental localizer generated by its aspherical projections. If factorizations exist, a $W$-map between intrinsic fibrations over a common base remains in $W$ after every base change: factor the base map into $W$ followed by an intrinsic fibration and apply two-out-of-three to the resulting square. Now factor the graph of any $f:X\to Y$ in $W$ as $X\xrightarrow{i}P\xrightarrow{q}X\times Y$, with $i\in W$ and $q\in\Fib_W$. Product projections belong to $\Fib_W$, so both $r=\operatorname{pr}_Xq$ and $p=\operatorname{pr}_Yq$ are intrinsic and in $W$. They are therefore universally in $W$. Pulling $p$ back to $Y/y$ shows $P/y$ is $W$-aspherical, hence $W_{\rm asph}$-aspherical. Absolute Theorem A gives $p\in W_{\rm asph}$, and likewise $r\in W_{\rm asph}$. From $ri=1$ and $pi=f$, it follows that $f\in W_{\rm asph}$. Thus $W=W_{\rm asph}$; Proposition 6.1.8 and Theorem 6.1.11 of \citet{RESCisinski2006} identify this with properness when $W$ is accessible.

Here is the required accessibility argument for the two rational localizers. Choose an uncountable regular cardinal $\kappa$ greater than $|\mathbb Q|$. For a weak arrow $u:A\to B$ start with any subfunctor of size less than $\kappa$. For rational homology, enlarge it by choosing finite chain witnesses for surjectivity and injectivity on every homology group. For rational local coefficients, first choose finite zigzags witnessing bijectivity on components, then do the same for the rational chain complexes of every based loop space of $\operatorname{Ex}^{\infty}NA\to\operatorname{Ex}^{\infty}NB$. These are the exact tests of D.5. Each simplex in the Kan replacement, each loop simplex, and each chain has finite support in the original nerves. There are fewer than $\kappa$ witnesses at each stage. Closing their arrows under composition and repeating countably often produces a weak subfunctor of size less than $\kappa$. The weak subfunctors so obtained form a $\kappa$-filtered system whose union is $u$: any fewer than $\kappa$ subfunctors have an upper bound by the same enlargement. The two tests commute with filtered colimits, and every fundamental localizer is closed under filtered colimits by \citet[Corollary 4.2.22]{RESCisinski2006}. Consequently representatives of the small weak subfunctors generate each localizer as a fundamental localizer. This proves the needed set-generation, without inferring it merely from closure under filtered colimits.

To see they are not proper, use Lemma 9.4.5 of \citet{RESCisinski2006}: a proper localizer contained in $W_r$ makes a weak equivalence with local Kan target a strong $r$-equivalence, in particular surjective on $\pi_{r+1}$. For $W_{\mathbb Q}$ take $K(\mathbb Z,1)\to K(\mathbb Q,1)$ and $r=0$. For $W_{D(\mathbb Q)}^{\mathrm{lc}}$ take $K(\mathbb Z,2)\to K(\mathbb Q,2)$ and $r=1$. These are equivalences for the indicated localizers, but their displayed homotopy-group maps are not onto. Their targets are local: the homotopy groups of maps into $K(\mathbb Q,j)$ are rational cohomology groups in degrees at most $j$, invariant under the corresponding equivalences. In the second case rational local coefficient equivalences imply ordinary rational homology equivalences by the universal-cover spectral sequence. This verifies the local-target hypothesis, and finishes the table.
\end{proof}

\subsection{Complete covering tests for ordinary integral and rational homology}
\begin{claim}{D.12}{Proved}{Trivial monodromy is necessary and sufficient for integral and rational homology}
Let $p:E\to B$ be a categorical covering, that is, a discrete fibration and discrete opfibration. For either $k=\mathbb Z$ or $k=\mathbb Q$, it belongs to $\Fib_{W_{Hk}}$ if and only if its monodromy on every fibre set is trivial. Equivalently, over each connected component $B_0$ it is isomorphic to a product $B_0\times S_0\to B_0$, where $S_0$ may be empty or infinite. In particular a nonempty connected covering of a connected base is intrinsic for either of these two homology localizers exactly when it is isomorphic over the base to its identity covering. This assertion is specific to these coefficients; D.8 gives the contrary mod-$p$ behavior.
\end{claim}
\begin{proof}
A categorical covering defines a set-valued local system on the groupoid completion of $B$. Conversely the unique source and target lifts reconstruct the covering from this local system. It is a product on a connected component precisely when every based loop acts trivially; choose paths from a basepoint to give a compatible trivialization. A homology equivalence with either indicated nonzero coefficient ring is a bijection on components and is a homology equivalence on each component. Products with a fixed discrete set, hence arbitrary disjoint unions of copies, preserve such equivalences. This proves sufficiency after every base change, including empty and infinite fibres.

For necessity choose a loop in $BB$ acting by a nonidentity permutation $P$ of the fibre set $S$. First let $k=\mathbb Q$. Pull the covering back to a circle representing this loop, and compare it with its pullback by the degree-$m$ map of the circle. Every nonzero degree is a rational homology equivalence. The components of these two coverings are the orbits of $P$ and $P^m$, respectively. If $P$ has a finite orbit of length $d>1$, take $m=d$; this orbit splits into $d$ distinct $P^m$-orbits. If it has an infinite orbit, take $m=2$; that orbit splits into two. The induced map on zeroth rational homology has a nonzero kernel. Thus the pullback is not a rational homology equivalence.

For $k=\mathbb Z$ we use a different witness. Fix $n\geq2$. In the universal cover of $S^1\vee S^n$, the Hurewicz isomorphism identifies $\pi_n$ with $\mathbb Z[t,t^{-1}]$, with $t$ acting by the circle deck transformation. Attach one $(n+1)$-cell along the class $2t-1$, and call the resulting CW space $X$. Collapsing the $n$-sphere to the basepoint and extending across the attached cell gives a map $u:X\to S^1$ with the original circle as section. The extension exists because $\pi_nS^1=0$. Its relative cellular chain complex has two terms, with differential the augmentation of $2t-1$, namely $1$. Hence the inclusion $S^1\to X$ and its retraction $u$ are integral homology equivalences.

Pull the covering over $S^1$ back along $u$. Relative to the lifted circle its two-term cellular chain complex has differential
\[
 2P-1:\mathbb Z[S]\longrightarrow\mathbb Z[S].
\]
Decompose $S$ into permutation orbits. On an orbit of length $d$, its cokernel is $\mathbb Z/(2^d-1)$: the relations $2e_{i+1}=e_i$ reduce to one generator and $(2^d-1)e_0=0$. It is nonzero for $d>1$. On an infinite orbit the cokernel is
$\mathbb Z[t,t^{-1}]/(2t-1)\cong\mathbb Z[1/2]$, also nonzero. Therefore the relative homology of the pulled pair is nonzero whenever $P$ is nonidentity. Since the lifted projection has the lifted circle as a section, it cannot be an integral homology equivalence.

These space-level witnesses supply actual small-category tests over the original $B$, rather than tests over a replacement base alone. Apply the slice realization construction in D.6 to each triangle over $BB$, and then the category-of-simplices functor followed by last-vertex evaluation. A covering has the homotopy-pullback property: its unique transport lifts give the homotopy-fibre comparison, or this follows by integrating its set-valued local system. Thus its strict categorical pullbacks realize the two coverings in the argument. The first categorical map is a $W_{Hk}$-equivalence and its pullback is not, proving necessity. The space $X$ is completely specified by the indicated attaching homotopy class; any representative gives the same relative chain complex, and the slice construction may use its full singular complex, so no unstated finiteness claim is involved.

The program \texttt{code/root/verify\_coverings.py} checks the integral matrices on cycles of lengths one through twelve by Smith normal form, verifies augmentation one, and records determinants $2^d-1$. This is a check of the displayed chain calculation; the argument above, including infinite orbits and every base change, proves the universal assertion.
\end{proof}
